\documentclass{article}
\usepackage[margin=0.75in]{geometry}
\usepackage{amsmath,amssymb,amsfonts}
\usepackage{amsthm}
\usepackage{graphicx}
\usepackage{hyperref}
\usepackage{enumitem}
\usepackage{physics}
\usepackage{bm}
\usepackage{siunitx}
\usepackage{booktabs}
\usepackage{caption}
\usepackage{subcaption}
\usepackage{cleveref}
\usepackage{xcolor}

%  define custom commands
\newcommand{\R}{\mathbb{R}}
\newcommand{\N}{\mathbb{N}}
\newcommand{\Z}{\mathbb{Z}}
\newcommand{\E}{\mathbb{E}}
\newcommand{\Symmetric}{\mathbb{S}}
\newcommand{\diff}{\mathrm{d}}
\newcommand{\D}{\mathrm{d}}
\newcommand{\cov}[1]{\mathrm{Cov}[#1 ]}
\let\P\relax
\newcommand{\P}[1]{\mathbb{P}\left[#1\right]}
\let\E\relax
\newcommand{\E}[1]{\mathbb{E}\left[#1\right]}
\newcommand{\textr}{\textcolor{red}}
\newcommand{\textb}{\textcolor{blue}}
\newcommand{\suchthat}{\mathrm{s.t.}\quad}
\newcommand{\normal}{\mathcal{N}}
\newcommand{\indicator}{\mathds{1}}
\newcommand{\chol}[1]{(#1)^{\frac{1}{2}}}
\newcommand{\maxeig}[1]{\lambda_{\max} \left\{#1 \right\}}
\newcommand{\blkdiag}[1]{\mathrm{blkdiag}(#1)}
\newcommand{\diag}[1]{\mathrm{diag}(#1)}
\DeclareMathOperator*{\minimize}{minimize}


\newtheorem{theorem}{Theorem}
\newtheorem{assumption}{Assumption}
\newtheorem{problem}{Problem}
\newtheorem{remark}{Remark}

\usepackage{cleveref}
\crefname{align}{Eq.}{Eqs.}
\crefname{equation}{Eq.}{Eqs.}
\crefname{figure}{Fig.}{Figs.}
\crefname{table}{Table}{Tables}
\crefname{theorem}{Theorem}{Theorems}
\crefname{definition}{Definition}{Definitions}
\crefname{lemma}{Lemma}{Lemmas}
\crefname{remark}{Remark}{Remarks}
\crefname{assumption}{Assumption}{Assumptions}
\crefname{proof}{Proof}{Proofs}
\crefname{algorithm}{Algorithm}{Algorithms}
\crefname{problem}{Problem}{Problems}
\crefname{proposition}{Proposition}{Propositions}
\crefname{corollary}{Corollary}{Corollaries}
\crefname{section}{Section}{Sections}


\allowdisplaybreaks[1] %allows equations/align to span multiple pages

\newcommand\blfootnote[1]{%
  \begingroup
  \renewcommand\thefootnote{}\footnote{#1}%
  \addtocounter{footnote}{-1}%
  \endgroup
}

\def\showChanges{0} % Change this value to 1 (red) or 0 (black) to switch highlighting condition

% Define the highlighting command
\newcommand{\highlight}[1]{%
    \ifnum\showChanges=1
        \textcolor{red}{#1}%
    \else
        #1%
    \fi
}
% highlighting command for math mode, to preserve spacing
\newcommand{\mathhl}[1]{
    \ifnum\showChanges=1
        \mathcolor{red}{#1}
    \else
        #1
    \fi
}
\newcommand{\needsfix}[1]{%
    \textcolor{blue}{#1}%
}

\usepackage{marginnote}
\usepackage{tcolorbox}
\setlength{\marginparwidth}{0.8in}
\newcommand{\margincomment}[1]{%
    \ifnum\showChanges=1
    \textcolor{blue}{\setlength{\baselineskip}{10pt}\footnotesize\raggedright\marginnote{#1}}%
    \else
    \fi
}

% set path for figures
\graphicspath{{../figures/}}

\title{Commented Content from Square Root QR Covariance Steering}
\author{Naoya Kumagai}

\begin{document}

\maketitle

\section{Keywords}
Covariance steering, square root methods, QR decomposition, chance constraints, sequential convex programming

\section{Full-size QR Decomposition}
\begin{lemma}[Existence of Full-size QR decomposition, Thm 3.3.3 in (cite)]
	Let \( M \in \mathbb{R}^{m \times n} \), \(m > n\). Then there exist \(Q \in \mathbb{R}^{m \times m}\) and \(\tilde{R} \in \mathbb{R}^{m \times n}\) such that \(Q\) is orthogonal, \( \tilde{R} = \begin{bmatrix}
	R \\ 0_{(m-n) \times n}
	\end{bmatrix}\), where \( R \in \mathbb{R}^{n \times n} \) is upper triangular, and \( M = QR \).
\end{lemma}

\subsubsection{Existence and Uniqueness}

\section{QR Decomposition Derivative}
In order to use this square root covariance propagation in an optimization framework, we utilize the derivative of the QR decomposition operation (i.e. the sensitivity of the $Q$ and $R$ matrices with respect to perturbations in the $M$ matrix). $\dd R$ is written as a function of \(Q\), \( R\), and  $\dd M$ \cite{deleeuwDifferentiatingQRDecomposition2023}:
\begin{equation}
	A + (\dd R) R^{-1} = Q^\top (\dd M) R^{-1}
\end{equation}
where \( A \) is a skew-symmetric matrix such that 
\begin{equation}
	\mathrm{tril}(A) = \mathrm{tril}(Q^\top (\dd M) R^{-1}).
\end{equation}
\( \mathrm{tril}(A) \) replaces the upper triangular part of \( A \) with zeros.
Organizing the above, we have:

\section{Implementation Notes}
Although $\bm{v}$ or $\bm{\mu}$ only appear in a convex form throughout the problem, the calculation of $\Delta J$ and $\Delta L$ requires their reference values at the first iteration.

For this problem, the short horizon length enables the block-based approach to also solve this problem with reasonable computational cost.

\section{Appendix}

A notable difference from typical trust region constraints is that the variable \( S_k \) is defined on a manifold. As described in \cite{kraislerIntrinsicSuccessiveConvexification2025a}, the Riemannian metric is used to define the trust region constraint. 

\subsection{Riemannian geometry on Cholesky spaces}
Let \( L \in \mathcal{L}_+^n \).
For \( X, Y \) in the tangent space of \( \mathcal{L}_+^n \) at \( L \), the Riemannian metric is given by \cite{linRiemannianGeometrySymmetric2019}:
\begin{equation}
	g_L(X,Y)
 = \sum_{i>j} X_{ij} Y_{ij} + \sum_{jj} X_{jj} Y_{jj} L_{jj}^{-2}
 \end{equation}

 Then, as described in \cite{kraislerIntrinsicSuccessiveConvexification2025a}, the trust region constraint is given by:
 \begin{equation} \label{eq:trust_region_constraint}
 	g_{\bar{S}_k}(S_k - \bar{S}_k, S_k - \bar{S}_k) \leq r^2
 \end{equation}
 where \( r > 0 \) is the trust region radius. This is a quadratic convex constraint.

 For \(L_k\), we simply impose a Euclidean trust region constraint:
 \begin{equation} \label{eq:trust_region_constraint_L}
 	\|L_k - \bar{L}_k\|_F \leq r_L
 \end{equation}

\begin{equation} \label{eq:trust_region_constraint_vectril}
	\mathrm{vectril}(S_k - \bar{S}_k) \leq r \quad k=0,\cdots, N
\end{equation}
where \( r > 0 \) is the trust region radius. 

\subsection{Scaling}
For problems where the covariance matrices have small numerical values, a numerical rescaling of the problem can improve numerical stability for the underlying conic solver. First, 
\begin{lemma}
	The QR decomposition $X=QR$ is invariant to the diagonal scaling of the matrix. In other words, for a diagonal matrix $D$, we have $XD = Q (RD)$.
\end{lemma}
\begin{proof}
	\(X=QR \implies XD = Q (RD)\). Since post-multiplication by $D$ does not change the rank of a matrix, from the uniqueness of the QR decomposition, the result follows.
\end{proof}

\subsection{Scaling for Covariance Steering}
From the previous lemma, a scaled version of \cref{eq:square_root_covariance_propagation} can be written as follows:
\begin{align*}
	\bar{S}_{k+1} = D S_{k+1} &= \mathrm{qr}( \begin{bmatrix} A_k S_k + B_k K_k S_k & G_k \end{bmatrix}^\top D)^\top \\
	&= \mathrm{qr}(\begin{bmatrix} D A_k D^{-1} D S_k + D B_k K_k D^{-1} D S_k & D G_k \end{bmatrix}^\top)^\top \\
	&= \mathrm{qr}(\begin{bmatrix} \bar{A}_k \bar{S}_k + \bar{B}_k \bar{K}_k \bar{S}_k & \bar{G}_k \end{bmatrix}^\top)^\top \\
	&= \mathrm{qr}(\begin{bmatrix} \bar{A}_k \bar{S}_k + \bar{B}_k \bar{L}_k & \bar{G}_k \end{bmatrix}^\top)^\top
\end{align*}
where $\bar{A}_k = D A_k D^{-1}$, $\bar{B}_k = D B_k$, $\bar{G}_k = D G_k$, $\bar{K}_k = K_k D^{-1}$, $\bar{L}_k = \bar{K}_k \bar{S}_k = K_k D^{-1} D S_k = K_k S_k = L_k$ and $\bar{S}_k = D S_k$.

The objective function can written in terms of the scaled variables as follows:
\begin{align*}
	J = \sum_{k=0}^{N-1} \mu_k^\top Q_k \mu_k + v_k^\top R_k v_k + \mathrm{trace} (\bar{Q}_k \bar{S}_k \bar{S}_k^\top + R_k \bar{L}_k \bar{L}_k^\top) \quad \text{(quadratic cost)}
\end{align*}
where $\bar{Q}_k = D^{-1} Q_k D^{-1}$. For the quantile 2-norm cost, the cost function is the same.

\section{Acknowledgments}
Acknowledge AI use
This material is based upon work supported by the Air Force Office of Scientific Research under award number FA9550-23-10512. N. Kumagai acknowledges support for his graduate studies from the Shigeta Education Fund.

\bibliographystyle{plain}
\bibliography{references}

\end{document}

